\chapter{Related work}
\label{ch:related_work}

As software development continues to evolve, it has seen the rise of virtual machines, containers, and various tools for automation.
This chapter aims to provide a consolidated review of the existing literature.

Studies by \citet{kolyasnikovExperimentalEvaluationVirtual2019} and \citet{felterUpdatedPerformanceComparison2015} report on the notable performance advantages of containers over VMs, especially concerning startup time and resource usage.
On a similar note, \citet{ionescuAlternativesRunningLinux2019} details the intricacies of both hypervisors and containers, and their performance implications.

Docker's efficiency, especially in the realm of microservices, is highlighted by \citet{potdarPerformanceEvaluationDocker2020}, which acknowledges its benefits for cloud deployments in terms of scalability and resource optimization.

\citet{merkelDockerLightweightLinux2014} discusses the lightweight nature of containers and their suitability across various software development phases.
The reproducibility challenges, especially in research domains, are addressed by \citet{hungGUIdockUsingDocker2016}, showcasing the potential of containers like Docker in ensuring consistent results.
Meanwhile, \citet{citoUsingDockerContainers2016} and \citet{chamberlainUsingDockerSupport2014} focus on Docker's role in aiding reproducible research.

\citet{gohringPracticalGuideContainerize2022} offers valuable insights into containerizing C applications.
\citet{koskinenContainersSoftwareDevelopment2019} presents a systematic mapping study that reveals a trend toward modularizing software systems using containers, with Docker emerging as the front-runner.

The criticality of continuous integration in modern agile setups is brought forth by \citet{agarwalContinuousIntegratedSoftware2018}, who advocates for automated testing and deployment to maintain software quality.
Dealing with large-scale C++ projects, \citet{gygiCppBuildLargeScaleAutomatic2021} presents the importance of automation in managing dependencies and ensuring reproducibility.
Additionally, \citet{randalIdealRealRevisiting2021} sheds light on the potential role of containers in simplifying software development complexities.

Integration of dockerized build environments, especially with tools like Visual Studio Code, is the central theme of the works of \citet{danilovDockerizedBuildEnvironments2021,danilovVSCodeDockerized2022}.
They discuss the benefits of such integration in providing a consistent development environment and addressing challenges faced by developers.
For another hands-on approach, \citet{luDevelopmentDockerContainers2018} offers guidance on C++ development using Docker within Visual Studio Code.

\section{Summary}

While the literature offers extensive insights into the merits of containerization, automation in software development, and the advantages of using development containers, a significant research gap remains regarding the effective integration of containers within long-standing or very big legacy software development projects and the ensuing benefits.